% =========================================================================
% Figura: Bifurcación de metodologías CADD (SBDD y LBDD) y convergencia en VS
% Inspirado en Chawla et al. (2024) y Macalino et al. (2015)
% Paleta de colores accesible Okabe-Ito definida en Configurations/01-Colors.sty
% =========================================================================

\begin{figure}[htbp]
\centering
\resizebox{0.80\linewidth}{!}{%
\begin{tikzpicture}[
    font=\sffamily,
    >=Stealth,
    x=1cm, y=1cm
]

% -------------------------------------------------------------------------
% 1. NODO RAÍZ SUPERIOR: CADD
% -------------------------------------------------------------------------
\node[
    draw=grayText!80,
    fill=cardBorder!40,
    rounded corners=5pt,
    line width=1.2pt,
    minimum width=4.6cm,
    minimum height=1.0cm,
    align=center,
    text=darkText
] (cadd) at (7.0, 5.95) {
    {\fontsize{12}{14}\selectfont \textbf{CADD}}\\[2pt]
    {\footnotesize\color{grayText} Diseño de Fármacos Asistido por Computadora}
};

% -------------------------------------------------------------------------
% 2. BIFURCACIÓN PRINCIPAL: SBDD Y LBDD
% -------------------------------------------------------------------------
% Rama SBDD (Izquierda - Azul)
\node[
    draw=figBlue,
    fill=figBlueTint,
    rounded corners=4pt,
    line width=1.2pt,
    minimum width=4.8cm,
    minimum height=0.9cm,
    align=center,
    text=darkText
] (sbdd) at (3.6, 4.15) {
    {\fontsize{11}{13}\selectfont \textbf{\color{figBlue} SBDD}}\\[1pt]
    {\footnotesize (\textit{Structure-Based Drug Design})}
};

% Rama LBDD (Derecha - Naranja)
\node[
    draw=figOrange,
    fill=figOrangeTint,
    rounded corners=4pt,
    line width=1.2pt,
    minimum width=4.8cm,
    minimum height=0.9cm,
    align=center,
    text=darkText
] (lbdd) at (10.4, 4.15) {
    {\fontsize{11}{13}\selectfont \textbf{\color{figOrange!90!black} LBDD}}\\[1pt]
    {\footnotesize (\textit{Ligand-Based Drug Design})}
};

% Conexión bifurcada de CADD a SBDD y LBDD con bifurcación equidistante
\path (cadd.south) -- (sbdd.north) coordinate[midway] (cadd_mid);
\coordinate (cadd_split) at (cadd.south |- cadd_mid);
\draw[line width=1.1pt, draw=grayText]
    (cadd.south) -- (cadd_split);
\draw[line width=1.1pt, draw=grayText, -{Stealth[length=5pt, width=4pt]}]
    (cadd_split) -| (sbdd.north);
\draw[line width=1.1pt, draw=grayText, -{Stealth[length=5pt, width=4pt]}]
    (cadd_split) -| (lbdd.north);

% -------------------------------------------------------------------------
% 3. CAJAS DE METODOLOGÍAS / HERRAMIENTAS
% -------------------------------------------------------------------------

% --- Sub-cajas SBDD (2 cajas: Docking y Diseño de novo) ---
\node[
    draw=figBlue!70,
    fill=white,
    rounded corners=3pt,
    line width=0.9pt,
    minimum width=2.1cm,
    minimum height=1.05cm,
    align=center,
    text=darkText,
    font=\footnotesize
] (box_docking) at (2.3, 2.20) {
    \textbf{Docking}\\[2pt]\textbf{molecular}
};

\node[
    draw=figBlue!70,
    fill=white,
    rounded corners=3pt,
    line width=0.9pt,
    minimum width=2.1cm,
    minimum height=1.05cm,
    align=center,
    text=darkText,
    font=\footnotesize
] (box_denovo) at (4.9, 2.20) {
    \textbf{Diseño}\\[2pt]\textbf{\textit{de novo}}
};

% Conexión de SBDD a sus 2 cajas
\path (sbdd.south) -- (box_docking.north) coordinate[pos=0.45] (sbdd_mid);
\coordinate (sbdd_split) at (sbdd.south |- sbdd_mid);
\draw[line width=0.9pt, draw=figBlue!80]
    (sbdd.south) -- (sbdd_split);
\draw[line width=0.9pt, draw=figBlue!80, -{Stealth[length=4.5pt, width=3.5pt]}]
    (sbdd_split) -| (box_docking.north);
\draw[line width=0.9pt, draw=figBlue!80, -{Stealth[length=4.5pt, width=3.5pt]}]
    (sbdd_split) -| (box_denovo.north);

% --- Sub-caja LBDD (1 caja: QSAR) ---
\node[
    draw=figOrange!70,
    fill=white,
    rounded corners=3pt,
    line width=0.9pt,
    minimum width=2.6cm,
    minimum height=1.05cm,
    align=center,
    text=darkText,
    font=\footnotesize
] (box_qsar) at (10.4, 2.20) {
    \textbf{QSAR}
};

% Conexión de LBDD a su caja
\draw[line width=0.9pt, draw=figOrange!80, -{Stealth[length=4.5pt, width=3.5pt]}]
    (lbdd.south) -- (box_qsar.north);

% -------------------------------------------------------------------------
% 4. NODO COMÚN DE CONVERGENCIA: CRIBADO VIRTUAL (VS)
% -------------------------------------------------------------------------
\node[
    draw=figGreen,
    fill=figGreenTint,
    rounded corners=5pt,
    line width=1.2pt,
    minimum width=5.6cm,
    minimum height=0.95cm,
    align=center,
    text=darkText
] (vs) at (7.0, 0.35) {
    {\fontsize{11}{13}\selectfont \textbf{\color{figGreen!80!black} Cribado Virtual (VS)}}
};

% Colectores hacia el Cribado Virtual con altura suficiente para las flechas
\coordinate (vs_in_sbdd) at ($(vs.north)+(-1.6,0)$);
\coordinate (vs_in_lbdd) at ($(vs.north)+(1.6,0)$);
\coordinate (vs_turn_level) at ($(vs.north)+(0, 0.40)$);

% Colector SBDD
\draw[line width=0.9pt, draw=figBlue!70] (box_docking.south) -- ++(0,-0.25) coordinate (b_dock);
\draw[line width=0.9pt, draw=figBlue!70] (box_denovo.south) -- ++(0,-0.25) coordinate (b_deno);
\draw[line width=0.9pt, draw=figBlue!70] (b_dock) -- (b_deno);
\coordinate (b_mid) at ($(b_dock)!0.5!(b_deno)$);
\draw[line width=1.1pt, draw=figBlue, -{Stealth[length=5pt, width=4pt]}]
    (b_mid) -- (b_mid |- vs_turn_level) -- (vs_in_sbdd |- vs_turn_level) -- (vs_in_sbdd);

% Colector LBDD
\draw[line width=1.1pt, draw=figOrange, -{Stealth[length=5pt, width=4pt]}]
    (box_qsar.south) -- (box_qsar.south |- vs_turn_level) -- (vs_in_lbdd |- vs_turn_level) -- (vs_in_lbdd);

\end{tikzpicture}%
}
\caption{Bifurcación de las estrategias de diseño de fármacos asistido por computadora (CADD) en métodos basados en estructura (SBDD) y basados en ligando (LBDD), convergiendo en el cribado virtual. Adaptado de~\cite{chawla2024computational, macalino2015role}.}
\label{fig:bifurcacion_cadd}
\end{figure}
